\HMTNarrativeHeading{Ce que conserve une transformation}

La théorie des opérateurs donne une forme calculable à la conservation de l'état et de ses registres. Une lecture partielle obtient une composante ; le complément recueille la contribution qui permet de restituer le bilan complet. La somme des carrés des normes de la composante lue et de son complément reproduit le carré de la norme de l'état initial. Les formules de récupération spécifient quelles données et quels opérateurs rendent possible l'inversion de la lecture.

Lors de l’itération du transport, chaque étape incorpore un complément et
prolonge une composante terminale. Le bilan fini conserve
explicitement ce terminal. Le passage à une profondeur arbitraire
étudie la convergence et détermine quand l’archive réunit toute la
norme. Cet ordre permet d’interpréter la mémoire comme une opération
de conservation et de récupération, sous des conditions vérifiables.

L’extension extérieure montre comment sont transportés les aires et les volumes
de la réalisation linéaire. Les contributions mixtes participent aux
produits extérieurs et à leurs bilans. Le changement de
représentation conserve les relations que ces objets mesurent.
La récupération maintient le lien avec l’orientation et les formes
utilisées par chaque preuve.

La monodromie et les boucles entre régimes géométriques ajoutent une dynamique déterminée. L'article compose des transformations elliptiques, paraboliques et hyperboliques, conserve leur mémoire et calcule leur réponse. Les symétries de l'action de transport produisent les charges variationnelles de ce domaine. L'invariance d'un observable est, quant à elle, formulée au moyen de sa relation avec le transport. Ces deux formes de conservation sont reliées en conservant leurs définitions particulières.

Les réalisations physiques retrouvent ensuite les sections d’action,
les lecteurs d’aire et les échelles préalablement engendrées. L’exemple
gaussien permet de calculer la pureté et le spectre d’une lecture
réduite à partir d’une préparation conjointe spécifiée. Ses
corrélations conservent l’information de l’état complet. Ce
résultat montre une conséquence opératoire de l’argument : le
transport détermine aussi bien ce qu’une observation reçoit que
l’information nécessaire pour comprendre sa relation avec la totalité.
